%------------------------------------------------------------------------------%
\documentclass[fleqn,utf8,aspectratio=169,12pt]{beamer}

\usepackage{calculo_varias_variaveis-1}

\title{Erro na Linearização}

%------------------------------------------------------------------------------%
\begin{document}

\addtitle

%------------------------------------------------------------------------------%
\section{Erro na Aproximação Linear}

%------------------------------------------------------------------------------%
\begin{frame}{Erro na Aproximação Linear}
  Se $f$ possui primeira e segunda derivadas parciais
  \pause
  sobre um conjunto aberto
  contendo um retângulo $R$ centrado em \((a, b)\)
  \pause
  e se $M$ é um
  limitante superior para os valores de \(\abs{f_{xx}}\),  \(\abs{f_{yy}}\)
  e \(\abs{f_{xy}}\) em $R$
  \pause
  \[
    \emph{\abs{f_{xx}} \leq M} \qquad
    \emph{\abs{f_{yy}} \leq M} \qquad
    \emph{\abs{f_{xy}} \leq M} \qquad
    \forall (x, y) \in R
  \]
  \pause
  então o erro na linearização
  \[
    E(x, y) = f(x, y) - L(x, y)
  \]
  \pause
  satisfaz
  \[
    \emph{\abs{E(x, y)} \leq \frac{M}{2} \big(\abs{x-a}+\abs{y-b}\big)^2}
  \]
\end{frame}

%------------------------------------------------------------------------------%
\section{Exemplos}

\include{I-04-Linearizacao-erro-ex1}

%------------------------------------------------------------------------------%
\section{Lista Mínima}

%------------------------------------------------------------------------------%
\begin{frame}{Lista Mínima}
  Cálculo Vol. 2 do Thomas 12\textsuperscript{a} ed. -- Seção 14.6
  \begin{enumerate}
    \item Estudar o texto da seção
    \item Resolver os exercícios: 33-38
  \end{enumerate}

  \vfill
  Atenção: A prova é baseada no livro, não nas apresentações
\end{frame}

%------------------------------------------------------------------------------%
\end{document}
%------------------------------------------------------------------------------%
